% Active Matter lab (Lab C) -- letter, two-column layout.
% Compile on Overleaf (pdfLaTeX): pdflatex main && bibtex main && pdflatex main && pdflatex main
% Requires references.bib and the two figure PDFs named below.
% Red \todo{...} entries need experimental details or unfinished results.
% Source-only CHECK comments identify analysis or figure changes still to make.
\documentclass[10pt,twocolumn]{article}

\usepackage[a4paper,margin=2.2cm]{geometry}
\usepackage{graphicx}
\usepackage{amsmath,amssymb,bm}
\usepackage{xcolor}
\usepackage[numbers,sort&compress]{natbib}
\usepackage[colorlinks=true,linkcolor=blue!50!black,citecolor=blue!50!black,urlcolor=blue!50!black]{hyperref}
\setlength{\columnsep}{7mm}
\setlength{\emergencystretch}{1em}
\setlength{\parskip}{0pt}
\setlength{\parindent}{1.2em}
\bibliographystyle{unsrtnat}

\newcommand{\todo}[1]{\textcolor{red}{\textbf{[TODO: #1]}}}
\newcommand{\nn}{\hat{\mathbf n}}
\newcommand{\rr}{\mathbf r}

\title{Active Matter - Lab C}
\author{Yael Yaari \and Netta Kaplan}
\date{November 15, 2026}

\begin{document}
% Full-width title and abstract above the two-column article body.
\twocolumn[{
\begin{@twocolumnfalse}
\maketitle

\begin{abstract}
Active matter is widespread in nature, spanning scales from microscopic swimming bacteria to the collective motion of animal groups. Unlike passive Brownian particles, active particles continuously consume energy to generate motion, giving rise to complex dynamics and emergent behaviors far from thermodynamic equilibrium.
Here, we use Hexbugs as a macroscopic model of active matter to explore how individual particle dynamics give rise to collective behavior and influence the motion of passive objects. At the single-particle level, we observe a transition from ballistic to diffusive motion, governed by the rotational diffusion of the particle. As the number of active particles increases, their interactions give rise to collective clustering and flocking behavior. Furthermore, we demonstrate that the motion of passive objects embedded in active matter is influenced by their geometry, highlighting the role of broken symmetry and chirality in determining their dynamics.
\end{abstract}
\vspace{1em}
\end{@twocolumnfalse}
}]

% ================= BEGIN TEMPORARY SHARED EQUATIONS =================
% These are Eqs. (3), (2), and (4), respectively, in Nettadraft.pdf.
% MERGE: move these labels to Netta's corresponding equations before deleting
% this block, or update the references below to her labels. Retain the rod-specific
% MSD and mean-displacement predictions in the passive-object section.
\section*{Shared Model Equations (Temporary)}
The active Brownian circle-swimmer model used in the single-particle section describes translation and orientation by~\cite{Balda2022,Kiechl2026}
\begin{equation}
  \begin{aligned}
    \dot{\mathbf r}(t)&=v_0\hat{\mathbf n}(t)+\sqrt{2D_t}\,\boldsymbol\xi(t),\\
    \dot\theta(t)&=\Omega+\sqrt{2D_r}\,\eta(t).
  \end{aligned}
  \label{eq:shared_abc}
\end{equation}
Here $\mathbf r$ is position, $v_0$ is propulsion speed, and $\hat{\mathbf n}=(\cos\theta,\sin\theta)$ is the unit orientation vector.
The coefficients $D_t$ and $D_r$ describe translational and rotational diffusion, and $\Omega$ is the systematic angular velocity.
The independent noises $\boldsymbol\xi$ and $\eta$ are zero-mean Gaussian white noises with correlations $\langle\xi_i(t)\xi_j(t')\rangle=\delta_{ij}\delta(t-t')$ and $\langle\eta(t)\eta(t')\rangle=\delta(t-t')$.
Setting $\Omega=0$ gives the nonchiral model; setting $D_t=0$ as well gives the simpler ABP equations in the single-particle section.

The mean-squared displacement is defined as
\begin{equation}
  \mathrm{MSD}(t)=\big\langle|\mathbf r(t_0+t)-\mathbf r(t_0)|^2\big\rangle_{t_0},
  \label{eq:shared_msd}
\end{equation}
where $t$ is the lag and $t_0$ is the starting time.
The orientational correlation predicted by the model is
\begin{equation}
  \begin{aligned}
    C(t)&=\big\langle\hat{\mathbf n}(t_0+t)\cdot\hat{\mathbf n}(t_0)\big\rangle_{t_0}\\
        &=e^{-D_r t}\cos(\Omega t).
  \end{aligned}
  \label{eq:shared_corr}
\end{equation}
The persistence time is $\tau_p=1/D_r$.
% ================== END TEMPORARY SHARED EQUATIONS ==================

\section*{Passive Object Dynamics}

% ---------------------------------------------------------------- Passive objects: theory and gears
Active particles can drive passive objects through repeated contact.
An asymmetric shape can convert this forcing into directed motion, an effect known as rectification, predicted~\cite{Angelani2009} and observed~\cite{DiLeonardo2010} for microscopic gears in bacterial baths.
We test whether object geometry similarly directs motion in a bath of Hexbugs.

Hexbugs can remain in concave corners and continue pushing, while flat faces allow sliding and escape~\cite{Balda2022,DiLeonardo2010}.
Unequal contacts on the two faces of a gear tooth can therefore produce a mean torque.
If the bath has no net turning bias, otherwise identical mirror-image gears should then rotate oppositely; individual Hexbugs, however, can turn with a preferred handedness~\cite{Kiechl2026}.
Symmetry permits this rotation but does not guarantee it.

All passive objects were free to translate and rotate on the surface, in a bath of 17 Hexbugs inside a $45\times27\,$cm rectangular arena.
We placed two gears of diameter $8\,$cm with opposite tooth orientations in the bath [Fig.~\ref{fig:spinner}(a)] and tracked the angle of a line joining two markers for $104\,$s; a linear fit within each $8\,$s window gave the angular velocity $\omega$, positive counter-clockwise.
\todo{Give the gears' material and mass. Confirm whether they are otherwise identical mirror images.}

\begin{figure*}[!t]
  \centering
  \includegraphics[width=\textwidth]{figures/fig2_gears.pdf}
  \caption{Rotation of freely moving gears with opposite tooth orientations in a bath of 17 Hexbugs. (a)~Gears A and B with their measured directions of rotation; scale bars, $5\,$cm. (b)~Accumulated rotation of gear~A (blue) and gear~B (orange) in turns, positive counter-clockwise. Dashed lines have slopes equal to the mean angular velocity $\bar\omega$ over the $8\,$s windows (labelled) and are offset from the data for clarity. Both gears are shown for $104\,$s, the duration of the gear-A recording, which ends when the gear was moved by hand.}
  \label{fig:spinner}
\end{figure*}
% CHECK Fig. 2: check angle unwrapping and sign convention of the gear angle.

The two gears rotate in opposite directions [Fig.~\ref{fig:spinner}(b)]: gear~A clockwise at $24\pm6\,^\circ$/s and gear~B counter-clockwise at $31\pm3\,^\circ$/s (mean $\pm$ standard error over 13 windows).
Gear~A turns clockwise in 12 of its 13 windows and gear~B counter-clockwise in all 13, although the speed varies between windows (standard deviations of $21$ and $10\,^\circ$/s), and the accumulated rotation grows at a nearly constant rate, parallel to the lines of slope $\bar\omega$.
This is consistent with tooth geometry setting the direction of rotation, as for asymmetric gears in bacterial baths~\cite{DiLeonardo2010}.

% ---------------------------------------------------------------- Passive objects: rod theory
Shape can also affect translation.
Hexbugs become trapped on the concave side of a passive wedge and push it~\cite{Balda2022}, and in simulations a curved rod is propelled by an active bath whereas a straight rod is not~\cite{Mallory2014}.
In a homogeneous, isotropic bath with no net handedness, symmetry forbids mean propulsion of a straight rod but allows a force along the bisector of an equal-arm L-shaped rod.
To test this, we tracked a straight rod and an L-shaped rod with $7\,$cm arms, of equal projected area, in the same bath [Fig.~\ref{fig:rods}(a)].
The tracked position is the centroid of the rod's silhouette, found together with the rod's orientation by matching a rotated template of the rod to each frame ($10$ frames/s).
\todo{Give the straight rod's dimensions and both rods' material and mass.}

We describe rod translation by the nonchiral limit ($\Omega=0$) of Eq.~\eqref{eq:shared_abc}, with no laboratory-fixed drift; $v_0$, $D_t$ and $D_r$ now describe the rod and are fitted separately from the single-Hexbug parameters.
For the L-shaped rod, $\hat{\mathbf n}$ points along the bisector from the corner towards the opening [Fig.~\ref{fig:rods}(a)], so corner-first motion corresponds to the propulsion term $-v_0\hat{\mathbf n}$, with $v_0\geq0$.
This white-noise description applies once correlations between successive pushes have decayed.
The MSD of Eq.~\eqref{eq:shared_msd} is then~\cite{Kiechl2026}
\begin{equation}
  \mathrm{MSD}(t)=4D_t t+\frac{2v_0^2}{D_r^2}
    \left(D_r t-1+e^{-D_r t}\right),
  \label{eq:rodmsd}
\end{equation}
which becomes $4D_t t+v_0^2t^2$ for $t\ll1/D_r$.
Writing $\Delta\mathbf r(t_0,t)=\mathbf r(t_0+t)-\mathbf r(t_0)$, the mean displacement along the initial bisector is~\cite{Balda2022}
\begin{equation}
  \begin{aligned}
    \Delta_\parallel(t)
      &\equiv\big\langle\Delta\mathbf r(t_0,t)\cdot\hat{\mathbf n}(t_0)\big\rangle_{t_0}\\
      &=-\frac{v_0}{D_r}\left(1-e^{-D_r t}\right),
  \end{aligned}
  \label{eq:dpar}
\end{equation}
while the mean perpendicular displacement, $\Delta_\perp\equiv\langle\Delta\mathbf r(t_0,t)\cdot[\hat{\mathbf z}\times\hat{\mathbf n}(t_0)]\rangle_{t_0}$, is zero.

% ---------------------------------------------------------------- Passive objects: rod results
% Analysis: stage2_code/stage2_pole_runs.py (one run per video). Excluded: 17bugs_l_4 (25.5 s usable).
% Usable durations: L 90, 45, 51, 39, 79, 58 s; straight 107, 96, 91, 54 s. Hand frames are excluded, but
% displacement pairs may span them (the hand never touched the rods).
Each rod was filmed in separate videos, with the L-shaped rod started at a different orientation in each: six videos of $39$--$90\,$s usable duration for the L-shaped rod and four of $54$--$107\,$s for the straight rod.
A seventh L-shaped-rod video, with only $26\,$s usable, was excluded because it gives too few starting times at the longest lag of $20\,$s.
Frames with a hand in view are excluded.
All quantities are computed within each video and averaged over videos with equal weight; quoted uncertainties are standard errors across videos.
$v_0$ and $D_t$ are fitted to the MSD with Eq.~\eqref{eq:rodmsd} over $1$--$20\,$s, $D_r$ is estimated from $C(t)=e^{-D_r t}$ [Eq.~\eqref{eq:shared_corr}], and MSD exponents are fitted over $1$--$15\,$s.
% Length scale: inner arena 45 x 27 cm -> 24.21 px/cm (sd 0.24 over 13 videos x 2 axes; stage2_code/stage2_calibration.py).
% Cross-checks: L arm 7 cm = 164 px (23.4 px/cm), Hexbug 4.5 cm = 107 px (23.8 px/cm).
Lengths are calibrated by the inner size of the arena ($24.2\,$px/cm); the $3\%$ calibration uncertainty is not included in the quoted errors.

\begin{figure*}[!t]
  \centering
  \includegraphics[width=\textwidth]{figures/fig3_rods_videos.pdf}
  \caption{Translation of an L-shaped and a straight rod in a bath of 17 Hexbugs. (a)~The L-shaped (top) and straight (bottom) rods; arrows show the axes $\hat{\mathbf n}$, from the corner towards the opening for the L-shaped rod and along the length for the straight rod; scale bars, $5\,$cm. (b)~Mean displacement along ($\Delta_\parallel$, filled) and perpendicular to ($\Delta_\perp$, open) the initial axis, for the L-shaped rod (orange, six videos) and the straight rod (blue, four videos). Points are means over videos, and coloured bands are standard errors across videos. The dashed curve is Eq.~\eqref{eq:dpar}, with $v_0$ fitted to the MSD and $D_r$ from $C(t)$, not to $\Delta_\parallel$; the grey band shows $v_0$ varied by its standard error. The dashed zero line is the prediction for $\Delta_\perp$ and for the straight rod. The straight rod has no preferred end, so the sign of its $\hat{\mathbf n}$ was drawn at random for each video.}
  \label{fig:rods}
\end{figure*}

The straight rod's MSD grows nearly linearly, with an exponent of $1.09\pm0.08$.
A fit with $v_0=0$ gives $D_t=0.12\pm0.03\,$cm$^2$/s, and both mean displacement components are consistent with zero [Fig.~\ref{fig:rods}(b)].
The L-shaped rod's MSD is two to four times larger over $1$--$20\,$s, with an exponent of $1.23\pm0.03$, and the fit gives $v_0=0.24\pm0.04\,$cm/s and $D_t=0.25\pm0.06\,$cm$^2$/s.
Its mean parallel displacement at $20\,$s, $-4.1\pm0.7\,$cm, agrees with the $-4.6\pm0.7\,$cm predicted by Eq.~\eqref{eq:dpar} from these fitted values, and its perpendicular displacement, $-0.5\pm0.9\,$cm, is consistent with zero.
The negative sign means corner-first motion, seen in all six videos ($-1.5$ to $-6.1\,$cm at $20\,$s).
The L-shaped rod turns slowly: $D_r=0.0055\pm0.0019\,$s$^{-1}$, so $\tau_p=1/D_r\approx180\,$s exceeds every video and is poorly determined.
Over the $20\,$s shown, however, Eq.~\eqref{eq:dpar} stays within $6\%$ of $-v_0t$, so the prediction hardly depends on $D_r$.

% ---------------------------------------------------------------- Passive objects: rod interpretation
A larger MSD alone does not establish propulsion: beads in a bacterial bath spread superdiffusively at short times without any asymmetry~\cite{WuLibchaber2000,Bechinger2016}.
Correlated pushes, possibly prolonged by retention at the concave corner~\cite{Balda2022}, could likewise enlarge the rods' MSD without a mean force; contact durations were not measured.
Even the straight rod, which symmetry forbids from being propelled, gives $v_0\approx0.07\,$cm/s when $v_0$ is left free in the MSD fit.

A laboratory-fixed drift, for example from a tilt of the table, could also mimic propulsion, because the rod barely rotates within a video.
Starting the L-shaped rod at different orientations separates the two: across the six videos the corner-first direction $-\hat{\mathbf n}$ points all around the circle, and in every video the rod's mean displacement lies within $50^\circ$ of it.
Adding a drift velocity $\mathbf V$ to Eq.~\eqref{eq:shared_abc}, the mean velocity of a video over $20\,$s intervals is $\mathbf u=-v_0\langle\hat{\mathbf n}\rangle+\mathbf V$, with $\langle\hat{\mathbf n}\rangle$ averaged over the same intervals.
% Drift fit (stage2_pole_runs.py, lag 20 s, cm/s): straight only V = (-0.031+-0.016, +0.005+-0.016);
% L only v0 = 0.196+-0.048, V = (-0.040+-0.042, -0.025+-0.045); L with V = 0: v0 = 0.192+-0.042;
% joint v0 = 0.191+-0.035, V = (-0.037+-0.025, -0.012+-0.026). Image axes are used as the laboratory frame;
% the two original recordings were filmed with the arena rotated ~8 deg in the image (neglected).
The straight rod, for which $v_0=0$, gives $|\mathbf V|=0.03\pm0.02\,$cm/s.
A joint fit of both rods gives $v_0=0.19\pm0.04\,$cm/s with $|\mathbf V|=0.04\pm0.03\,$cm/s, so the drift is about a fifth of the propulsion speed, and $v_0$ from the mean displacements agrees with the MSD value within the uncertainties.
A laboratory-fixed drift therefore does not explain the L-shaped rod's motion, which follows the rod's orientation, as expected from shape-driven propulsion.
The speed varies between videos (fitted $v_0$ from $0.06$ to $0.31\,$cm/s), and the origin of this spread, for example uneven Hexbug density or the rod's distance from the walls, was not studied.


\bibliography{references}
\end{document}